% GENERATED FILE. Edit canonical lessons, then run npm run generate:books.
% canonical-source-hash: fnv1a64:49bd6cf6121929e7
% canonical-lessons: GU-C286-guno, GU-C286-sipai, GU-C286-sanka, GU-C286-khub, GU-C286-chupavvu

\chapter{Calling the Police}
\label{ch:gu-calling-the-police}
\hlchaptermodality{286}

\begin{chapteropening}
\textbf{By the end of this chapter:} \emph{I can report a crime and say whom I suspect.}

Complete the last lesson of chapter 286: I can report a crime and say whom I suspect.
\end{chapteropening}

\section[\texorpdfstring{\gu{ગુનો} (guno)}{ગુનો (guno)}]{\gu{ગુનો} (guno) — a crime}

\label{lesson:GU-C286-guno}



\begin{quote}
Before the new one: say the Gujarati for to catch, then the Gujarati for people.
\end{quote}

\subsection*{You'll want to know: \gu{ગુનો}}
\textbf{\gu{ગુનો}} — \emph{guno} — \textquotedblleft{}a crime\textquotedblright{}.

\begin{quote}
\gu{ચોરી} \gu{એક} \gu{ગુનો} \gu{છે}.
\end{quote}

\begin{quote}
\emph{corī ek guno che.}
\end{quote}

\begin{quote}
Theft is a crime.
\end{quote}

A masculine noun.

\begin{cousinweb}
\textbf{\gu{ગુનો}} (\emph{guno}) is usually traced to Persian \emph{gunāh}, \textquotedblleft{}a sin, a fault\textquotedblright{}; in Gujarati it is an offence against the law. It is masculine: \textbf{\gu{મોટો} \gu{ગુનો}}, a serious crime.
\end{cousinweb}

\begin{grammarlens}[title={What you've built}]
One more word for telling what happened.
\end{grammarlens}

\subsection*{Your turn}
\begin{itemize}
\raggedright
  \item \emph{Say it:} \emph{guno}
  \item \emph{Say it:} \emph{guno}, once more
  \item \emph{Say it:} \emph{corī ek guno che.} — Theft is a crime.
  \item \emph{Recall:} say the Gujarati for to catch, then the Gujarati for people, then say \emph{guno} again
\end{itemize}

\begin{tcolorbox}[breakable,colback=teal!4,colframe=teal!35!black,title={Before you move on}]
What does \gu{ગુનો} mean? (\textquotedblleft{}A crime\textquotedblright{}.) Say it once more.
\end{tcolorbox}
\section[\texorpdfstring{\gu{સિપાઈ} (sipāī)}{સિપાઈ (sipāī)}]{\gu{સિપાઈ} (sipāī) — a constable}

\label{lesson:GU-C286-sipai}



\begin{quote}
Before the new one: say the Gujarati for people, then the Gujarati for a crime.
\end{quote}

\subsection*{You'll want to know: \gu{સિપાઈ}}
\textbf{\gu{સિપાઈ}} — \emph{sipāī} — \textquotedblleft{}a constable\textquotedblright{}.

\begin{quote}
\gu{જો} \gu{ત્યાં} \gu{એક} \gu{સિપાઈ} \gu{હોત} \gu{તો} \gu{તે} \gu{ચોરને} \gu{પકડત}.
\end{quote}

\begin{quote}
\emph{jo tyā̃ ek sipāī hot to te corne pakḍat.}
\end{quote}

\begin{quote}
If a constable had been there, he would have caught the thief.
\end{quote}

A masculine noun.

\begin{cousinweb}
\textbf{\gu{સિપાઈ}} (\emph{sipāī}) is usually traced to Persian \emph{sipāhī}, \textquotedblleft{}a soldier\textquotedblright{}, the word English took as \emph{sepoy}. In Gujarati it is the constable on the beat.
\end{cousinweb}

\begin{grammarlens}[title={What you've built}]
One more word for telling what happened.
\end{grammarlens}

\subsection*{Your turn}
\begin{itemize}
\raggedright
  \item \emph{Say it:} \emph{sipāī}
  \item \emph{Say it:} \emph{sipāī}, once more
  \item \emph{Say it:} \emph{jo tyā̃ ek sipāī hot to te corne pakḍat.} — If a constable had been there, he would have caught the thief.
  \item \emph{Recall:} say the Gujarati for people, then the Gujarati for a crime, then say \emph{sipāī} again
\end{itemize}

\begin{tcolorbox}[breakable,colback=teal!4,colframe=teal!35!black,title={Before you move on}]
What does \gu{સિપાઈ} mean? (\textquotedblleft{}A constable\textquotedblright{}.) Say it once more.
\end{tcolorbox}
\section[\texorpdfstring{\gu{શંકા} (śankā)}{શંકા (śankā)}]{\gu{શંકા} (śankā) — a suspicion}

\label{lesson:GU-C286-sanka}



\begin{quote}
Before the new one: say the Gujarati for a crime, then the Gujarati for a constable.
\end{quote}

\subsection*{You'll want to know: \gu{શંકા}}
\textbf{\gu{શંકા}} — \emph{śankā} — \textquotedblleft{}a suspicion\textquotedblright{}.

\begin{quote}
\gu{મીરાને} \gu{ચોકીદાર} \gu{પર} \gu{શંકા} \gu{હતી}.
\end{quote}

\begin{quote}
\emph{Mīrāne cokīdār par śankā hatī.}
\end{quote}

\begin{quote}
Mira suspected the watchman.
\end{quote}

A feminine noun.

\begin{cousinweb}
\textbf{\gu{શંકા}} (\emph{śankā}) is a tatsama, Sanskrit \emph{śaṅkā}, \textquotedblleft{}fear, distrust\textquotedblright{} (CDIAL 12258). The person suspected takes \textbf{\gu{પર}}: \textbf{\gu{ચોકીદાર} \gu{પર} \gu{શંકા}}, a suspicion about the watchman.
\end{cousinweb}

\begin{grammarlens}[title={What you've built}]
One more word for telling what happened.
\end{grammarlens}

\subsection*{Your turn}
\begin{itemize}
\raggedright
  \item \emph{Say it:} \emph{śankā}
  \item \emph{Say it:} \emph{śankā}, once more
  \item \emph{Say it:} \emph{Mīrāne cokīdār par śankā hatī.} — Mira suspected the watchman.
  \item \emph{Recall:} say the Gujarati for a crime, then the Gujarati for a constable, then say \emph{śankā} again
\end{itemize}

\begin{tcolorbox}[breakable,colback=teal!4,colframe=teal!35!black,title={Before you move on}]
What does \gu{શંકા} mean? (\textquotedblleft{}A suspicion\textquotedblright{}.) Say it once more.
\end{tcolorbox}
\section[\texorpdfstring{\gu{ખૂબ} (khūb)}{ખૂબ (khūb)}]{\gu{ખૂબ} (khūb) — very; a great deal}

\label{lesson:GU-C286-khub}



\begin{quote}
Before the new one: say the Gujarati for a constable, then the Gujarati for a suspicion.
\end{quote}

\subsection*{You'll want to know: \gu{ખૂબ}}
\textbf{\gu{ખૂબ}} — \emph{khūb} — \textquotedblleft{}very; a great deal\textquotedblright{}.

\begin{quote}
\gu{ચોરી} \gu{પછી} \gu{મીરા} \gu{ખૂબ} \gu{ડરી} \gu{ગઈ}.
\end{quote}

\begin{quote}
\emph{corī pachī Mīrā khūb ḍarī gaī.}
\end{quote}

\begin{quote}
After the theft Mira was badly frightened.
\end{quote}

\begin{cousinweb}
\textbf{\gu{ખૂબ}} (\emph{khūb}) is usually traced to Persian \emph{khūb}, \textquotedblleft{}good, fine\textquotedblright{}. In Gujarati it makes a word stronger: \textbf{\gu{ખૂબ} \gu{સારું}}, very good; \textbf{\gu{ખૂબ} \gu{ડરી} \gu{ગઈ}}, got very frightened.
\end{cousinweb}

\begin{grammarlens}[title={What you've built}]
One more word for telling what happened.
\end{grammarlens}

\subsection*{Your turn}
\begin{itemize}
\raggedright
  \item \emph{Say it:} \emph{khūb}
  \item \emph{Say it:} \emph{khūb}, once more
  \item \emph{Say it:} \emph{corī pachī Mīrā khūb ḍarī gaī.} — After the theft Mira was badly frightened.
  \item \emph{Recall:} say the Gujarati for a constable, then the Gujarati for a suspicion, then say \emph{khūb} again
\end{itemize}

\begin{tcolorbox}[breakable,colback=teal!4,colframe=teal!35!black,title={Before you move on}]
What does \gu{ખૂબ} mean? (\textquotedblleft{}Very; a great deal\textquotedblright{}.) Say it once more.
\end{tcolorbox}
\section[\texorpdfstring{\gu{છુપાવવું} (chupāvvũ)}{છુપાવવું (chupāvvũ)}]{\gu{છુપાવવું} (chupāvvũ) — to hide (something)}

\label{lesson:GU-C286-chupavvu}



\begin{quote}
Before the new one: say the Gujarati for a suspicion, then the Gujarati for very.
\end{quote}

\subsection*{You'll want to know: \gu{છુપાવવું}}
\textbf{\gu{છુપાવવું}} — \emph{chupāvvũ} — \textquotedblleft{}to hide (something)\textquotedblright{}.

\begin{quote}
\gu{ચોરે} \gu{પૈસા} \gu{પથ્થર} \gu{નીચે} \gu{છુપાવ્યા}.
\end{quote}

\begin{quote}
\emph{core paisā paththar nīce chupāvyā.}
\end{quote}

\begin{quote}
The thief hid the money under a stone.
\end{quote}

Past, with \gu{તેણે} (he or she …): \gu{તેણે} \gu{પૈસા} \gu{છુપાવ્યા}, \emph{teṇe paisā chupāvyā}.

\begin{cousinweb}
\textbf{\gu{છુપાવવું}} (\emph{chupāvvũ}) is the causative of \emph{chupāvũ}, to be hidden, which CDIAL lists under \emph{*chupp}, \textquotedblleft{}to hide\textquotedblright{}, a form it reconstructs (CDIAL 5058). You hide a thing, and the verb agrees with it: \textbf{\gu{પૈસા} \gu{છુપાવ્યા}}.
\end{cousinweb}

\begin{grammarlens}[title={What you've built}]
One more word for telling what happened.
\end{grammarlens}

\subsection*{Your turn}
\begin{itemize}
\raggedright
  \item \emph{Say it:} \emph{chupāvvũ}
  \item \emph{Say it:} \emph{chupāvvũ}, once more
  \item \emph{Say it:} \emph{core paisā paththar nīce chupāvyā.} — The thief hid the money under a stone.
  \item \emph{Recall:} say the Gujarati for a suspicion, then the Gujarati for very, then say \emph{chupāvvũ} again
\end{itemize}

\begin{tcolorbox}[breakable,colback=teal!4,colframe=teal!35!black,title={Before you move on}]
What does \gu{છુપાવવું} mean? (\textquotedblleft{}To hide (something)\textquotedblright{}.) Say it once more.
\end{tcolorbox}
